
\section{\revision{Proof of Theorem \ref{Thm_withou_A1}}} 
\label{app:withoutA1}

\revision{
In this appendix, we describe the changes required in the proofs of
Theorems~\ref{Thm_Lim_SD_Conv}-\ref{Thm_PB} to establish
Theorem~\ref{Thm_withou_A1}. Note that we no longer make Assumption
{\bf A}.1, but simply assume that eager job sizes are exponentially
distributed.

Without Assumption {\bf A}.1, one can not describe the pre-limit
system using a one-dimensional Markov process. However with
exponential service times for both eager and tolerant customers, one
can describe the system evolution via a two-dimensional Markov chain
$(X_\tau (t), X_\i (t)),$ capturing the number of tolerant and eager
customers in the system. Note that this Markov chain undergoes
transitions due to $\i$ as well as $\t$ customers. As noted before,
when the occupancy of the tolerant changes, a new eager sub-policy
comes into effect immediately, potentially during an ongoing eager
busy period. This new sub-policy might result in scheduling
rearrangements, based on the number of eager customers in the system.

The main difference without Assumption {\bf A}.1 is that once the
$\tau$ occupancy changes, the first $\i$-busy cycle can start with
(random) $c$ number of customers, where
$0 \le c \le {\cal K} := \lfloor 1/c_{min} \rfloor < \infty. $ Let the
duration of this partial first $\i$-busy cycle be represented by
$\B_{j}^{\mue, c} $, when the system is operating under $j$-th
sub-policy; basically, this is the time period that starts immediately
after the $\t$-system changes to state $j$, and lasts till the
$\i$-system becomes empty (or till the $\t$-state changes again,
whichever happens first).\footnote{Recall that \emph{complete}
  $\i$-busy cycles start when the $\i$-system becomes empty and last
  till the end of the next $\i$-busy period.
  % that started with a single $\i$-customer.
} Aside from this first partial $\i$-busy-cycle, subsequent $\i$-busy
cycles are referred to using the same notations as before.
% For the quantities that correspond to the
% case in which the first busy cycle starts with $0$
% $\i$-customers\footnote{ Recall here that these $\i$-busy cycles start
%   when the $\i$-system becomes empty and last till the end of the
%   $\i$-busy period, that started with one $\i$-customer.  } (as with
% assumption {\bf A}.1) will be referred by same notations as before.
% We use special notations only for the case when the first
% \textcolor{green}{$\i$-}busy cycle starts with $c$ $\i$-customers.

While the two-dimensional Markovian description
$(X_\tau (t), X_\i (t)),$ with a random number~$c$ of eager customers
present at times of $\t$-transitions is more challenging to analyse,
we show below that the first dimension of this process can be upper
and lower bounded by one-dimensional Markov processes (as in
\cite{ANOR}). These bounding processes are described next.
% Thus, the main challenge in analysing the two dimensional process is
% because of $c$, one can still analyze the original process by using
% one dimensional processes that upper and lower bound the original
% process as in \cite{ANOR}.


% The stationary distribution exists and the time limit exists because
% one can dominate it ... and then equivalent of Lemma 2 shows
% stationry of dominant chain which in turn implies staitonary of
% original chains ? Or we may have to go through Maximum theorem or
% some other functional analysis tools.. Or do we have some queueing
% results such that the staitonary distreibution occurs for general
% queueing systems, need to use them.}

{\bf Upper and Lower bounding systems:} Now we construct two systems
for each $\mue$, which sandwich the actual $\t$-system pre-limit. In
both the bounding systems we consider that the eager system,
immediately after each tolerant change, starts with maximum possible
(partial) eager busy cycle: generate one busy cycle $\B_{j}^{\mue, c}$
starting with each $c$ and with coupled eager arrivals, job-sizes
etc., for the given sub-policy (say sub-policy~$j$) and let
$\B_{j}^{\mue, *}:= \max_{0 \le c \le {\cal K}} \B_{j}^{\mue, c}$ be
duration for which the first (hypothetical) eager busy cycle lasts in
the two systems. It is immediate that this max-busy cycle upper bounds
those starting with any $c$-number (as in original system), when they
are served with the same sub-policy.
% , using certain coupling arguments; every such busy period
% eventually hits $c$ number of eager customers and that residual busy
% cycle can be coupled with the one starting with $c$-number and this
% is possible due to exponential random quantities.  {\color{blue} We
% may generate one busy cycle starting with each $c$ and with coupled
% eager arrivals, job-sizes etc., and consider that both the bounding
% systems start with maximum of these ${\cal K}$ number of busy
% cycles.}

In the upper bounding system, immediately after the $\t$-occupancy
changes to~$j$, the tolerant queue receives zero service for duration
$\B_{j}^{\mue, *}$. After this period, service of the tolerant queue
in subsequent $\i$ busy cycles continues exactly as in the original
system.
% the $\i$-system starts with corresponding max-busy cycle, during
% which the $\t$-customer(s) are not served; the service of the
% $\t$-customer(s) starts with the second $\i$-busy cycle and
% continues exactly as in the original system.
The lower bounding system is almost similar to the upper bounding
system, except that the tolerant queue is served at the maximum
possible rate (1, given the unit server speed assumed) over the first
hypothetical partial busy cycle $\B_{j}^{\mue, *}.$
% , the $\t$-customers are served with maximum possible rate during
% the first $\i$ max-busy cycle.
We then use appropriate coupling rules as in \cite{ANOR}: \\

\noindent $\bullet$ Tolerant inter-arrival times and job-sizes are
always coupled in all three systems. However, the eager quantities are
coupled in the original and any bounding system, only when the two
systems are in the same $\t$-state, as explained below.\\

\noindent $\bullet$ The original and the two bounding systems start
with the same initial system state;\\

\noindent $\bullet$ the number of $\t$-customers in the original
system is less than or equal to that in upper system, and greater than
or equal to that in the lower system, as time progresses (because of
the service differences for the tolerant queue in the first partial
eager busy-cycles across the three systems);\\

  
\noindent $\bullet$ if at any time point the upper system
meets\footnote{The tolerant number in upper system can become strictly
  bigger than that in original systems, but further increase may be
  slower in upper system due to Markov policies, and then at some
  tolerant departure the $\t$-number in upper system can equal that in
  the original system.}  the original system, we couple the two
systems by using eager random variables (inter-arrival times, service
times, etc.) that define the original system in that $\t$-cycle to
define the upper system $\t$-cycle;\\
  
\noindent $\bullet$ specifically, we start with the maximum partial
busy cycle in upper system corresponding to the sub-policy dictated by
the new tolerant state, after coupling (also) the eager inter-arrival
times and job sizes; observe that under this construction, the first
$\i$-busy cycle in the original system (in the $\t$-cycles in which
upper system meets the original system) is smaller than or equal to
the corresponding maximum partial busy cycle in upper system.\\

These details are similar to those in \cite{ANOR}, where one can also
find more details. Similar coupling ideas can be used to construct the
lower bounding system.

It is easy to observe that the tolerant queue evolution in the two
bounding systems can be modeled as M/G/1 systems with state dependent
birth-death transition probabilities, as the system considered in the
paper with Assumption~{\bf A.}1.  That is, they can be represented by
one-dimensional Markov processes.  All the results can be extended to
the two bounding systems and we would have equivalents of Theorems
\ref{Thm_Lim_SD_Conv}-\ref{Thm_tau} and Lemma~\ref{lemma:TauStability}
for the two systems, if we can show that the impact of the first
maximum partial busy cycle at the start of every tolerant change
vanishes as $\mue \ra \infty$ under the SFJ limit. We show below that
this is the case.

{\bf Convergence of the bounding systems:} This proof goes through
along similar lines as with {\bf A}.1, except for few important
differences. We only mention the differences here.  We first show that
the max busy cycles $\B_{j}^{\mue, *}$ converge to zero, as
$\mu_\i \to \infty$, the moments are upper bounded by a common uniform
bound, etc. For extending Lemmas
\ref{lemma:bp_upper_bound}-\ref{lemma:mginftyBPbound}, we observe that
the max busy cycles can be uniformly upper-bounded (uniformly over
sub-policy $j$) by that in the M/G/$\infty$ queue of
Lemma~\ref{lemma:mginftyBPbound}, which starts with ${\cal K}$
customers.
  
For extending Lemma~\ref{lemma:uniform}, the transition probabilities
also depend upon max busy cycles. Towards this, first define
$q_j^{\mue, *}$ to represent the backward transition probability of
tolerant birth-death process when one starts it with corresponding max
busy cycle.  Then one can lower bound the probability $q_j^{\mue, *}$
using
$$\bar{q}^*_j := \prob{\S_{j,1}^{\mue,*} +\sum_{k = 2}^{N-1} \S_{j,k}^{\mue} > B_{\t}} = (1-r_*) {\bar q}_j + r_*, \  \    r_* := \prob { \S_{j,1}^{\mue,*}  > B_\tau},   $$  
with ${\bar q}_j$ as defined previously and where $\S_{j,1}^{\mue,*} $
equals the service available to the tolerant queue during the first
$\i$ max busy cycle $\B_{j}^{\mue, *}$ (which exactly equals 0 for
upper system and $\B_{j}^{\mue, *}$ in the lower system for any
$\mue$). It is clear that $r_*$ converges to zero (uniformly in $j$)
and the remaining proof goes through, as before.  The proof for the
upper bound of Lemma~\ref{lemma:uniform} can be handled similarly.
Observe here that Lemmas
\ref{Lemma_Intermediate_result}-\ref{Lemma_Intermediate_result_1} are
used by Lemma~\ref{lemma:uniform} only for the subsequent $\i$-busy
cycles, which are same as that with Assumption~{\bf A}.1, and hence
would not require any changes.  The proof of
Lemma~\ref{Lemma_Intermediate_result_2} holds also for the two
bounding systems, because the convergence in Lemma~\ref{lemma:uniform}
is ensured and same is the case with Lemma~\ref{lemma:TauStability}.

Thus, by the extended Lemma~\ref{lemma:TauStability} there exists
${\bar \mu}$ such that the upper and lower bounding systems are stable
for all $\mue \ge {\bar \mu}$.  Further the upper and lower stationary
distributions converge to the same quantity as $\mue \to \infty$ (by
extended versions of Theorems~\ref{Thm_Lim_SD_Conv}-\ref{Thm_tau}).

{\bf Existence of stationary distribution of original system:} We
derive the existence using embedded chains, the discrete time Markov
chains obtained by considering values of the continuous time Markov
process, immediately after the transition epochs.

By extended Lemma~\ref{lemma:TauStability} there exists a ${\bar \mu}$
such that the upper bounding system is stable for all
$\mue \ge {\bar \mu}$. For all such $\mue$, the embedded tolerant
chain in the upper system visits the~$0$ state infinitely often (with
probability one) and within integrable number of epochs (by positive
recurrence and existence of stationary distribution). This implies (by
the bounding of the tolerant occupancy between the original and upper
system) that the embedded tolerant chain in the original system also
visits $0$ state infinitely often (with probability one) and within an
integrable number of epochs. This in turn implies the original
(two-dimensional) embedded Markov chain visits the set
$ {\cal O}:= \{ (0, c), \mbox{ for some } c \le {\cal K} \}$
infinitely often (w.p.1) and within integrable number of epochs.  By
irreducibility, this implies the same for the $(0,0)$ state, as every
time the chain visits any state of the form $(0,c),$ it has a positive
probability of visiting state $(0, 0)$ (uniformly bounded away from 0
across all $c$) and hence visits state $(0,0)$ after at most a
geometric number of visits to the set~${\cal O}.$ This establishes the
existence of a unique stationary distribution
% (uniqueness by aperiodicity)
for the original two-dimensional embedded Markov chain, which in turn
implies the existence of stationary distribution for the two
dimensional continuous time Markov process; observe that the
transition rates between all the state pairs can be uniformly upper
bounded by a finite value for any given $\mue.$
 
By regular renewal arguments, this implies almost sure existence of
the time limits:
 \begin{eqnarray}
\label{Eqn_SD_marginal}
\lim_{T \to \infty}  \frac{1}{T} \int_0^T  \indc_{ \{X_\tau(t) = j \}}  =  \pi_j^\mue, \mbox{ for any } j, 
\end{eqnarray}
where $\pi_j^\mue$ is the marginal of the stationary distribution
corresponding to tolerant queue in original (two-dimensional) system.

{\bf Convergence of stationary distribution as $\mue\to \infty$:} The
limit of the $\t$-performance of the two bounding systems, as
$\mue \to \infty$ is exactly the same, because by extended
Lemma~\ref{lemma:uniform} both the sets of transition probabilities
converge to the same limit transition probabilities uniformly.
Furthermore, time averages corresponding to the $\t$-occupancy in the
original system can be bounded by the corresponding time averages in
the upper system (occupancy denoted by $X_\t^U (\cdot)$) and the lower
system (occupancy denoted by $X_\t^L (\cdot)$) as follows:
%by bounding arguments
%($X_\t^U (.)$, $X_\t^L(.)$ correspond respectively to upper and lower
%systems),
$$
\lim_{T \to \infty} \frac{1}{T} \int_0^T \indc_{ \{X^{L}_\tau(t) \le j
  \}} \le \lim_{T \to \infty} \frac{1}{T} \int_0^T \indc_{ \{X_\tau(t)
  \le j \}} \le \lim_{T \to \infty} \frac{1}{T} \int_0^T \indc_{
  \{X^U_\tau(t) \le j \}} \mbox{ for any } j,
$$
and hence the marginal stationary distribution (corresponding to
tolerant occupancy) of the original system, converges to the common
limit of the two bounding systems. This implies the statement related
to the tolerant performance in Theorem \ref{Thm_Lim_SD_Conv}, because
(for each $j$):
\begin{eqnarray*}
\sum_{i\le j} \pi_i  = \lim_{\mue \to \infty } \lim_{T \to \infty}  \frac{1}{T} \int_0^T  \indc_{ \{X^{L}_\tau(t) \le  j \}}  \le  \lim_{\mue \to \infty }    \lim_{T \to \infty}  \frac{1}{T} \int_0^T  \indc_{ \{X_\tau(t) \le  j \}}  \hspace{30mm}  \\
 \le
 \lim_{\mue \to \infty }   \lim_{T \to \infty}  \frac{1}{T} \int_0^T  \indc_{ \{X^U_\tau(t) \le  j \}} =  \sum_{i \le j} \pi_i, \  \    \mbox{ w.p.1. }
\end{eqnarray*}
Using similar bounding arguments, and common limit of the two bounding
systems, the statement of Theorem \ref{Thm_tau} also gets extended.


{\bf Changes required for results related to eager performance:} There
are no changes in the main proof of Theorem \ref{Thm_PB}, we only need
to discuss the changes required in the proof of
Theorem~\ref{Thm_PB_for_one_subpolicy}.

For extending the proof of Theorem \ref{Thm_PB_for_one_subpolicy}, one
needs to construct additional sample paths for coupling arguments.
One first needs to construct ${\cal K}+1$ additional sample paths of
$\i$-busy cycles, with $\mue = 1$, one for each sub-policy~$j,$ and
one for each $0 \le c \le {\cal K}$.  Using these ${\cal K}+1$ sample
paths, we need to construct a $({\cal K}+2)$th sample path of max busy
cycles, one for each $j$.  Note again that each such max busy cycle
upper bounds the corresponding busy cycle starting with $c$ number of
$\i$-customers, for any $0 \le c \le {\cal K}$. Whenever the
$\t$-state changes, the first $\i$-busy cycle is coupled using one of
the above ${\cal K}+1$ sample paths, depending upon the random $c$,
that represents the number of eager customers at the time of the
$\t$-transition. Each such first partial busy cycle is is upper
bounded by the corresponding max busy cycle, irrespective of the
random~$c$; this ensures the required coupling (and uniform upper bounding across various $c$) across various $\mue$,
even without Assumption~{\bf A}.1.

The process $\Psi_{j} (t)$ for each $j$-sub-policy now includes the
starting as well as the ending partial busy cycles.  But all the
starting partial $\i$-busy cycles can be upper bounded by the
corresponding max busy cycles and one can show the influence of them
to converge to zero, in the same manner as is done with the ending
busy cycles; this is possible because we have already extended the
proofs of the required uniform bounds (converging to zero or second moments bounded
as required) on all the required moments of the max busy cycles, as
already discussed in the initial parts of this section. All the bounds
that were shown to converge to zero, should now be made up of the
corresponding max busy cycles, to handle the part of $\Psi_{j} (t)$
coming from the starting partial $\i$-busy cycles.  The same is true
for the bounds and arguments considered in the proof of
Lemma~\ref{Lemma_Psij}.

If there any drops due to change in eager sub-policy at tolerant transition epochs, they can be upper bounded by  $N_{\partial_j} (t) $ and can be handled as in original proof. 
%Finally, we will not have $N_{\partial_j} (t) $, the losses due to
%drops at $\t$-transitions (without {\bf A}.1). 
\eop

}
